\documentclass[11pt,a4paper]{article}

\usepackage[T1]{fontenc}
\usepackage[utf8]{inputenc}
\usepackage[cyr]{aeguill}
\usepackage[francais]{babel}

% %\renewcommand{\baselinestretch}{1.5}  % line spacing : 1,5
% \setlength{\oddsidemargin}{2,5cm}	% Marge gauche sur pages impaires
% \setlength{\evensidemargin}{2,5cm}	% Marge gauche sur pages paires
% \setlength{\textwidth}{16cm}	% Largeur de la zone de texte (17cm)
% \setlength{\topmargin}{2,5cm}	% Pas de marge en haut

% Les packages
%\usepackage[french]{babel}
\usepackage{fancybox}
\usepackage{graphics}
\usepackage{color}
\usepackage{latexsym}
\usepackage{amsmath,amsthm,amssymb}
\usepackage[plain]{fullpage} 
\usepackage{verbatim}
\usepackage{times,epsfig}
\usepackage{listings}
\usepackage{multirow}
\usepackage{amsmath,amsthm,amssymb,amscd,txfonts,bbm} % RAJOUT ICI !!!
\usepackage{graphics,epsfig} % RAJOUT ICI !!!
%\usepackage{algorithmic}
\usepackage[final]{pdfpages} 
\usepackage{fancyvrb}
\usepackage{enumerate}
\usepackage{listingsutf8}
\usepackage{comment}


\newenvironment{qv}
{\quote\Verbatim}
{\endquote\endVerbatim}

\usepackage{lastpage}

\newcommand{\terme}[1]{\texttt{#1}}
\newcommand{\termSet}[1]{\{\texttt{#1}\}}
\newcommand{\guill}[1]{<<\,#1\,>>}
% Macro pour l'inclusion de figure  
\newcommand{\fig}[2]{%
  \begin{figure*}[hbt]
   \begin{center}
      \includegraphics[width=6cm]{#1}
  \caption{\label{#1}  #2}
  \end{center}
 \end{figure*}
}
\newcommand{\figTex}[2]{%
  \begin{figure*}[hbt]
 \centerline{\input{#1.pstex_t}}
  \caption{\label{#1}  #2}
 \end{figure*}
}

% Macro commentaires
\definecolor{turquoise}{rgb}{0.20 ,0.60, 0.60}
\definecolor{vert}{rgb}{0.40,  0.60 ,0.00}
\definecolor{violet}{rgb}{0.80 , 0.00 , 0.80}
\definecolor{bleu}{rgb}{0.20,  0.20,  0.80}
\definecolor{rose}{rgb}{1.00,  0.20,  0.40}
\def\remVG#1{\textcolor{vert}{\texttt{Virginie: #1}}}
\def\comVG#1{\begin{footnotesize}\textcolor{vert}{\texttt{Precision de Virginie: #1}}\end{footnotesize}}

\lstset{basicstyle=\footnotesize,showstringspaces=false,language=XML,breaklines=false,columns=fullflexible,frame=shadowbox, captionpos=b}

\newenvironment{solution}[0]{
~\\ \color{red} \textbf{Solution : }\color{black}}{%
}
%\excludecomment{solution}


\parindent=0pt           
                             
\usepackage{url}

\newtheorem{Exercice}{Definition}

\lstset{language=XML,basicstyle=\footnotesize,
numberstyle=\footnotesize, 
breaklines=true
showspaces=false,         
showstringspaces=false,
showtabs=false,
frame=single,
tabsize=4}   


\begin{document}

\title{Examen NFE204 - Session 2}
\date{14 avril 2015}

\maketitle

\begin{center}
{\bf \Large Consignes}\\
\begin{tabular}{| p{15cm} |}
\hline
\\
Dur\'ee de l'examen : 3h00\\
Documents non autoris\'es;  Calculatrice autoris\'ee\\

Les exercices sont indépendants les uns des autres. Merci de soigner votre écriture.\\
Ce sujet contient 2 pages.\\


\textit{Important}\,: nous n'évaluons pas les réponses par la syntaxe. Ne vous
inquiétez pas d'utiliser un point-virgule à la place d'un point d'exclamation ou
autres détails mineurs
\\
\hline
\end{tabular}
\end{center}

\section*{Première partie : recherche d'information (8 pts)}


\begin{tabular}{lp{15cm}}
     $d_1$&Le loup et les trois petits cochons.\\
     $d_2$&Spider Cochon, Spider Cochon, il peut marcher au plafond. Est-ce qu'il
         peut faire une toile? Bien sûr que non, c'est un cochon.\\
     $d_3$&Un loup a mangé un mouton, les autres moutons sont restés dans la bergerie.\\
     $d_4$&Il y a trois moutons dans le pré, et un autre dans la gueule du loup.\\
     $d_5$&L'histoire extraordinaire des trois petits loups et du grand méchant cochon.\\ 
\end{tabular}

\smallskip

On va se limiter au vocabulaire \texttt{loup, mouton, cochon}.

\begin{enumerate}
  \item Donnez la matrice d'incidence \emph{booléenne} (seulement 0 ou 1)
      avec les termes en ligne (NB: on suppose
     une phase préalable de normalisation qui élimine les pluriels, majuscules, etc.)

\begin{solution}
\begin{center}
\begin{tabular}{l|c|c|c|c|c}
           & \textbf{d1} & \textbf{d2} & \textbf{d3} &\textbf{d4}&\textbf{d5} \\ \hline  
   loup    & 1           & 0           & 1           & 1         & 1\\
   mouton  & 0           & 0           & 1           & 1         & 0\\
   cochon  & 1           & 1           & 0           & 0         & 1\\\hline
   \end{tabular}
\end{center}
\end{solution}
    
  \item Expliquez par quelle technique, avec cette matrice d'incidence, on peut 
     répondre à la requête booléenne "loup et cochon mais pas mouton".
     Quel est le résultat?
     
\begin{solution}
On fait un \texttt{AND} logique sur les vecteurs "loup" et "cochon"
et sur la négation du vecteur "mouton".
\end{solution}

   \item Maintenant donnez une matrice d'incidence contenant les tf, 
      et un tableau donnant les idf (sans le $\log$), pour les trois termes précédents.

\begin{solution}
\begin{center}
\begin{tabular}{l|c|c|c|c|c}
            & \textbf{d1} & \textbf{d2} & \textbf{d3} &\textbf{d4}&\textbf{d5} \\ \hline  
   loup  5/4   & 1           & 0           & 1           & 1         & 1\\
   mouton 5/2  & 0           & 0           & 2           & 1         & 0\\
   cochon 5/3  & 1           & 3           & 0           & 0         & 1\\\hline
   \end{tabular}
\end{center}
\end{solution}

  \item Donner les résultats classés par similarité cosinus basée sur les tf 
  (on ignore l'idf) pour les requêtes suivantes.
  
  \begin{itemize}
    \item cochon;
    \item loup et mouton;
    \item loup et cochon.
  \end{itemize}
\begin{solution}
Les normes
\begin{itemize}
  \item $||d_1|| = \sqrt{1+ 1}$; $||d_2|| = \sqrt{3^2}$; $||d_3|| = \sqrt{1+4}$; 
    $||d_4|| = \sqrt{1+1}$, $||d_5|| = \sqrt{1+1}$  
\end{itemize}

Les cosinus (requête non normalisée).

\begin{itemize}
  \item cochon: d1 : $\frac{1}{\sqrt{2}}$;  
       d2 : $\frac{3}{3}$; 
       d5 : $\frac{1}{\sqrt{2}}$; 
       
       d2 est premier: il ne parle que de cochons. d1 et d5
       sont à égalité.
  
  \item loup et mouton: 
  
    d1 : $\frac{1}{\sqrt{2}}$; 
    d3 : $\frac{1+2}{\sqrt{5}}$;
    d4 :  $\frac{1+1}{\sqrt{2}}$  
    d5 : $\frac{1}{\sqrt{2}}$; 
    
    d4 est premier: il a le même équilibre entre les deux termes que la requête.
    d3 est plus orienté mouton que loup, les deux derniers ne parlent
    pas de mouton.
    
  \item loup et cochon
  
    d1 : $\frac{1+1}{\sqrt{2}}$; 
    d2 : $\frac{3}{3}$;
    d3 :  $\frac{1}{\sqrt{5}}$  
    d4 :  $\frac{1}{\sqrt{2}}$  
    d5 : $\frac{1+1}{\sqrt{2}}$; 
    
  d1 et d5 sont en tête à égalité, devant d2, d4 et d3.
 \end{itemize}

\end{solution}  

   \item On ajoute le document $d_6$: "Shaun le mouton: une nuit
   de cochon".
   
   Quel est son score pour la requête "loup et mouton", 
   quel autre document a le même score et qu'est-ce qui change
   si on prend en compte l'idf?
   
\begin{solution}
d6 obtient le même score que d1, mais l'emporte si on considère
l'idf de "mouton" qui est plus élevé  (6/3 en comptant d6) que celui de "loup"
(6/4 en comptant d6).
\end{solution}

   \item Pour la requête "loup et cochon" et les documents d1 et d5, 
   qu'est-ce qui change si on ne met pas de restriction 
   sur le vocabulaire (tous les mots sont indexés)?
   
\begin{solution}
Dans ce cas le plus petit document (d1) l'emporte car sa norme est
plus petite. Intuititvement: il parle plus \emph{spécifiquement} de
loup et de cochon que d5.
\end{solution}   
    
  \end{enumerate}

\section*{Seconde partie : Pig et MapReduce (6 pts)}

Une organisation terroriste, le Spectre, envisage de commettre un attentat
dans une station de métro. Heureusement, le MI5 dispose d'une base de données d'échanges
téléphoniques et ses experts ont mis au point un décryptage qui identifie
la probabilité qu'un message provienne du Spectre d'une part, et fasse
référence à une station de métro d'autre part.

Après décryptage, les messages obtenus ont la forme suivante:

\begin{verbatim}
{
  "id": "x1970897",
  "émetteur": "Joe Shark", 
  "contenu": "Hmm hmm hmmm",
  "probaSpectre": 0.6,
  "station": "Covent Garden"
}
\end{verbatim}

Il y en a des milliards: vous avez quelques heures pour trouver la solution.

\medskip
\textbf{Questions}.

\begin{enumerate}
  \item Spécifier les deux fonctions du programme MapReduce qui
      va identifier la station-cible la plus probable. Ce programme doit sélectionner
     les messages émis par le Spectre avec une probabilité de plus de 70\%,  et produire, 
     pour chaque station le nombre de tels messages où elle est mentionnée.

\begin{solution}
function map(message)
{
if (message.probaSpectre > 70) {
emit (message.station, 1)
}
}

function reduce(station, groupe)
{
return <station, sum(groupe)>;
}
\end{solution}


  \item Quel est le programme Pig qui exprime ce traitement MapReduce?
  
  \item On veut connaître les 5 mots les plus couramment employés par 
   chaque émetteur dans le contenu de ses messages. Expliquez, avec la formalisation 
      de votre choix, comment obtenir cette information (vous avez le droit 
      d'utiliser un opérateur de tri).
\end{enumerate}


\begin{solution}
\begin{enumerate}
 \item 
\end{enumerate}
\end{solution}

\section*{Troisième partie : questions de cours (6 pts)}


\begin{enumerate}

  \item En recherche d'information, que signifient les termes "faux positifs"
      et "faux négatifs"?


\item Je soumets une requête $t_1, t_2, \cdots, t_n$. 
    Quel est le poids de chaque terme dans le vecteur représentant cette requête? 
    La normalisation de ce vecteur est elle importante pour le classement (justifier)?
    
  \item Donnez trois bonnes raisons de choisir un système relationnel plutôt
      qu'un système NoSQL pour gérer vos données.

  \item Donnez trois bonnes raisons de choisir un système NoSQL plutôt
      qu'un système relationnel pour gérer vos données.
  
  \item Rappeler la règle du quorum 
      (majorité des votants) en cas de partitionnement de réseau, et justifiez-là.

  \item Dans un traitement MapReduce, peut-on toujours se contenter d'un seul
     \textit{reducer}? Avantages? Inconvénients?
      
\end{enumerate}

\end{document}
